\documentclass[aspectratio=169]{beamer} % Alterado para 16:9 (padrão moderno)
\setbeamerfont{headline}{size=\footnotesize}

\mode<presentation>{
    \usetheme{Montpellier}
    \setbeamercovered{transparent}
    \usecolortheme{beaver} % Sugestão: Beaver tem tons de vermelho/cinza que funcionam bem, ou volte para lsc se preferir
}

\usepackage[english,brazil]{babel}
\usepackage[utf8]{inputenc}
\usepackage{fontawesome}
\usepackage{graphicx}
\usepackage{listings}
\usepackage{url}
\usepackage{colortbl}
\usepackage{caption}
\usepackage{tikz}
\usetikzlibrary{positioning,arrows.meta,fit,backgrounds,shapes.multipart,shapes, arrows, positioning, shadow, calc}

\usepackage{ragged2e} % Para justificar texto se necessário

% --- Configuração de Cores e Listings ---
\definecolor{codegreen}{rgb}{0,0.6,0}
\definecolor{codegray}{rgb}{0.5,0.5,0.5}
\definecolor{codepurple}{rgb}{0.58,0,0.82}
\definecolor{backcolour}{rgb}{0.95,0.95,0.92}

\newcommand{\reflexao}[1]{\begin{alertblock}{Reflexão}#1\end{alertblock}}
\newcommand{\key}[1]{\underline{#1}}
% --- Dados do Título ---
\title[Mineração de Dados]{Mineração de Dados: Conceitos e Aplicações}
\subtitle{}
\author[Elton Sarmanho]{
    Prof. Elton Sarmanho\inst{1}\\
    \footnotesize \href{mailto:eltonss@ufpa.br}{eltonss@ufpa.br}
}
\institute[UFPA]{
    \inst{1} Faculdade de Sistemas de Informação - UFPA Campus Cametá
}
\date{\today}

% --- Configuração de Diagramas TikZ ---
\tikzstyle{block} = [rectangle, draw, fill=blue!20, text width=5em, text centered, rounded corners, minimum height=4em]
\tikzstyle{line} = [draw, -latex']
\tikzstyle{cloud} = [draw, ellipse, fill=red!20, node distance=3cm, minimum height=2em]
\tikzstyle{decision} = [diamond, draw, fill=green!20, text width=4.5em, text badly centered, inner sep=0pt]

\begin{document}

% Slide de Título
\begin{frame}
    \titlepage
    \begin{center}
        \tiny \faGithub \ \href{https://github.com/eltonsarmanho}{github.com/eltonsarmanho} \quad
        \faLinkedinSquare \ \href{https://www.linkedin.com/in/elton-sarmanho-836553185/}{Elton Sarmanho}
    \end{center}
\end{frame}

% Roteiro
\begin{frame}{Roteiro da Aula}
    \tableofcontents
\end{frame}

% ==============================================================================
\section{Conceitos Fundamentais}
% ==============================================================================

\subsection{Titulo da subseção}
\begin{frame}{Nome do Slide}
   
\end{frame}



\begin{frame}{Referências Bibliográficas}
    \begin{thebibliography}{99}
        \bibitem{heuser} HEUSER, C. A. \textit{Projeto de Banco de Dados}. 6ª ed. Bookman, 2009.
        \bibitem{silber} SILBERSCHATZ, A. et al. \textit{Sistema de Banco de Dados}. Campus.
        \bibitem{date} DATE, C. J. \textit{Introdução a Sistemas de Banco de Dados}. Campus.
        \bibitem{navathe} ELMASRI, R.; NAVATHE, S. \textit{Sistemas de Banco de Dados}. Pearson.
    \end{thebibliography}
\end{frame}

\end{document}
